\section{Widening}

Nei domini astratti con catene ascendenti infinite, l'iterazione standard
per il calcolo del punto fisso potrebbe non terminare.
L'operatore di widening ($\nabla$) è una tecnica fondamentale per garantire
la convergenza dell'analisi in un numero finito di passi, forzando una
stabilizzazione della sequenza di iterati.

\subsection{Definizione dell'operatore di widening}

\begin{definizione}{Operatore di widening}{widening_operator}
Un operatore binario $\nabla: A \times A \to A$ è un operatore di widening
su un dominio astratto $A$ (che è un join-semilattice) se soddisfa due
proprietà fondamentali:

\begin{itemize}
    \item \textbf{Upper bound (Sovra-approssimazione)}:
    \[
    \forall x,y \in A : x \sqsubseteq x \nabla y \quad \text{e} \quad y \sqsubseteq x \nabla y
    \]
    Il risultato deve essere un upper bound di entrambi gli operandi (quindi
    anche della loro unione $x \sqcup y$).

    \item \textbf{Convergenza}:
    Per ogni catena ascendente $(y_i)_{i \in \mathbb{N}}$ in $A$, la
    sequenza $(x_i)_{i \in \mathbb{N}}$ definita come:
    \[
    x_0 = y_0, \quad x_{i+1} = x_i \nabla y_{i+1}
    \]
    si stabilizza in un tempo finito.
    Ovvero $\exists k \ge 0$ tale che $x_{k+1} = x_k$.
\end{itemize}
\end{definizione}

\begin{esempio}{Widening su intervalli}{interval_widening}
Nel dominio degli intervalli, l'operatore di widening (\Cref{def:widening_operator})
è progettato per proiettare all'infinito gli estremi che risultano instabili
(che crescono) tra un'iterazione e l'altra.

Dati due intervalli $I_1 = [a, b]$ e $I_2 = [c, d]$ tali che $I_1
\sqsubseteq I_2$, il widening è definito come:
\[
[a, b] \nabla [c, d] =
\left[
    \begin{cases} a & \text{se } a \le c \\ -\infty & \text{altrimenti} \end{cases},
    \begin{cases} b & \text{se } b \ge d \\ +\infty & \text{altrimenti} \end{cases}
\right]
\]
Se un estremo è stabile (es. $a \le c$ per il limite inferiore), viene
mantenuto.
Se invece l'estremo si sta ``allargando'' ($a > c$), il widening lo forza
immediatamente a $-\infty$ (o $+\infty$ per il limite superiore) per garantire
che smetta di cambiare.

Proprietà:
\begin{itemize}
    \item \textbf{Non commutativo}:
    $[2,3] \nabla [1,4] = [-\infty, +\infty] \neq [1,4] \nabla [2,3]$.
    \item \textbf{Non monotono}: Può accadere che pur avendo input più
    precisi, il risultato del widening sia meno preciso.
\end{itemize}

\end{esempio}

\subsection{Strategie di applicazione}

L'applicazione indiscriminata del widening a tutte le equazioni porterebbe a
una grossolana sovra-approssimazione.
Secondo Bourdoncle, è sufficiente applicare il widening solo in un insieme di
punti che interrompono ogni ciclo nel grafo delle dipendenze (feedback vertex
set), tipicamente le intestazioni dei cicli (loop headers).

Il sistema di equazioni astratto diventa quindi:
\[
x_i =
\begin{cases}
    x_i \nabla f_i^\#(x_1, \dots, x_m) & \text{se } i \in WP \text{ (Widening Points)} \\
    x_i \sqcup f_i^\#(x_1, \dots, x_m) & \text{altrimenti}
\end{cases}
\]

\section{Narrowing}

L'utilizzo del widening garantisce la convergenza ma introduce una perdita di
precisione significativa (spesso introducendo valori $\pm \infty$).
Il risultato ottenuto è un post-punto fisso ($postfp(f)$), tale che
$F^\#(x) \sqsubseteq x$.
L'operatore di \textbf{narrowing} ($\triangle$) permette di raffinare questo
risultato recuperando precisione.

\begin{definizione}{Operatore di Narrowing}{narrowing_def}
Un operatore binario $\triangle: A \times A \to A$ è un operatore di narrowing
se soddisfa le seguenti proprietà:

\begin{itemize}
    \item \textbf{Raffinamento limitato}:
    \[
    \forall x,y \in A : x \sqcap y \sqsubseteq x \triangle y \sqsubseteq x
    \]
    Il risultato deve essere contenuto nel primo operando $x$ (migliorando la
    precisione) ma deve comunque contenere l'intersezione reale $x \sqcap y$
    (per mantenere la soundness).

    \item \textbf{Convergenza}:
    Per ogni sequenza $(y_i)_{i \in \mathbb{N}}$, la catena $(x_i)_{i \in \mathbb{N}}$ definita come:
    \[
    x_0 = y_0, \quad x_{i+1} = x_i \triangle y_{i+1}
    \]
    si stabilizza in un tempo finito.
\end{itemize}
\end{definizione}

\begin{esempio}{Narrowing su Intervalli}{interval_narrowing}
Nel dominio degli intervalli, il narrowing permette di recuperare estremi
finiti partendo da estremi infiniti generati dal widening, ma solo se il
nuovo valore calcolato è un miglioramento valido.

\[
[a, b] \triangle [c, d] =
\left[
    \begin{cases} c & \text{se } a = -\infty \\ a & \text{altrimenti} \end{cases},
    \begin{cases} d & \text{se } b = +\infty \\ b & \text{altrimenti} \end{cases}
\right]
\]
La logica è la seguente:
\begin{itemize}
    \item Se il vecchio estremo era infinito ($-\infty$ o $+\infty$) e il
    nuovo valore proposto è finito, accettiamo il nuovo valore (abbiamo
    recuperato precisione).
    \item Se il vecchio estremo era già finito, lo manteniamo inalterato
    anche se il nuovo valore è ``migliore'' (più stretto).
    Questo è cruciale per garantire la terminazione ed evitare cicli infiniti
    di raffinamento.
\end{itemize}
\end{esempio}

\subsection{Fasi dell'Analisi}

L'analisi statica completa con domini infiniti si svolge in due fasi sequenziali:

\begin{enumerate}
    \item \textbf{Ascending Phase (Fase Ascendente)}:
    Si calcola la sequenza partendo da $\bot$ e utilizzando il widening
    ($\nabla$) nei punti di widening $WP$. Questa fase termina quando si
    raggiunge un post-punto fisso $x^\nabla$.
    \[ x_0 = \bot, \quad x_{i+1} = x_i \nabla F^\#(x_i) \]

    \item \textbf{Descending Phase (Fase Discendente)}:
    Si parte dal risultato della fase precedente $x^\nabla$ e si utilizza
    il narrowing ($\triangle$) per raffinare la soluzione.
    \[ y_0 = x^\nabla, \quad y_{i+1} = y_i \triangle F^\#(y_i) \]
    Il limite di questa sequenza $x^\triangle$ è un'approssimazione più
    precisa del least fixpoint reale.
\end{enumerate}

\begin{nota}{Widening e narrowing}{}
    
    \begin{itemize}
        \item \textbf{Complementarietà}: Widening e Narrowing sono due facce
        della stessa medaglia.
        Il primo ``spinge'' verso l'alto per garantire di coprire tutte le
        soluzioni (anche a costo di esagerare), il secondo ``tira'' verso il
        basso per eliminare l'eccesso.
        \item \textbf{Correttezza vs Precisione}: Il widening sacrifica la
        precisione per la terminazione.
        Il narrowing tenta di recuperare la precisione persa senza sacrificare
        nuovamente la terminazione.
        \item \textbf{Non monotonicità}: È interessante notare come il
        widening non sia monotono. Questo è controintuitivo rispetto alla
        maggior parte degli operatori sui reticoli, ma è necessario per
        ``saltare'' oltre i punti di accumulazione infiniti.
    \end{itemize}
\end{nota}